\chapter{Preliminary Analysis}
\label{ch:analysis}

\section{Tools Used}

The tools are described by category. The exact models and hosting options will
be chosen during implementation.

\begin{itemize}
  \item \textbf{Attack generation (red side):} open-source red-teaming tools and
  public jailbreak and prompt-injection datasets, used to create and vary attack
  prompts.
  \item \textbf{Target model:} an open-source LLM, run locally or through a
  cloud-hosted inference API. The final choice will be made during
  implementation.
  \item \textbf{Defense (blue side):} an open-source guardrail /
  safety-classifier toolkit to detect and block unsafe responses.
  \item \textbf{Reasoning / judge:} a reasoning model that is prompted to run the
  evidence-first and hypothesis-first checks.
  \item \textbf{Orchestration:} a Python framework that connects the agents, the
  reasoning engine, and the memory.
  \item \textbf{Storage:} a PostgreSQL database holds the auditable memory; logs
  are stored as JSON/JSONB records.
\end{itemize}

\section{Data Sources}

PurpleAgents uses three kinds of data. It is important to understand where each
one comes from.

\begin{itemize}
  \item \textbf{Hypotheses (attack categories):} a prepared list built by
  combining several public sources --- the OWASP Top 10 for LLM Applications, the
  OWASP Agentic Top 10, MITRE ATLAS, and open jailbreak and prompt-injection
  datasets. This list is defined once and is used by the hypothesis-first path.
  \item \textbf{Attack prompts:} drawn from open jailbreak and injection
  datasets, and then extended by the red agent, which generates new variations.
  \item \textbf{Benign prompts:} a small set of ordinary, harmless requests,
  taken from public instruction datasets. These are used to check that new fixes
  do not block normal use, and to measure the false-flag rate.
  \item \textbf{Evidence:} this is not downloaded. It is created live every time
  an attack runs, and consists of the prompt plus the model's actual response.
  This is the data that the reasoning engine works on.
\end{itemize}

\noindent\textbf{Note:} the project does not depend on a single fixed dataset.
Most of its data is generated during testing.

\section{Evaluation Metrics}

Table~\ref{tab:metrics} defines the metrics that will be used to evaluate
PurpleAgents. Here, an attack is \emph{successful} when the verified conclusion
says the target produced unsafe output and the blue agent did not block it. All
values will be measured during implementation; \textbf{no values are reported
in this report}, and any target levels mentioned are \emph{expected} only.

\begin{table}[H]
\centering
\caption{Evaluation metrics (all values to be measured; expectations only)}
\label{tab:metrics}
\small
\begin{tabularx}{\textwidth}{@{}P{0.19\textwidth} Y Y@{}}
\toprule
\textbf{Metric} & \textbf{Meaning} & \textbf{How it will be measured} \\
\midrule
Attack Success Rate (ASR) & Fraction of attacks that got through the defense. Lower is better. & $\mathrm{ASR}_r$ = successful attacks in round $r$ $\div$ total attacks in round $r$. \\
Detection Rate (DR) & Fraction of attacks that were caught (catch-rate). Higher is better. & $\mathrm{DR}_r = 1 - \mathrm{ASR}_r$, computed per round. \\
Improvement across rounds & How much the defense improved through learning. & $\Delta\mathrm{DR} = \mathrm{DR}_r - \mathrm{DR}_1$; expected to be positive, though not every round need improve. \\
Generalization & Whether improvements carry over to attacks the system never used for fixing. & DR on the held-out attack set, measured before and after the learning rounds, compared with DR on the attacks used for fixing. \\
False-flag rate & How often the defense wrongly blocks harmless requests. Lower is better. & Benign prompts blocked $\div$ total benign prompts, measured after each round to show that fixes do not harm normal use. \\
Disagreement rate & How often the two reasoning paths disagree and a case is sent to a human. & Cases where $A \neq B$ $\div$ total cases; a very high value would mean the judge prompts need improvement. \\
\bottomrule
\end{tabularx}
\end{table}

A good outcome is a detection rate that rises across rounds, holds up on the
held-out set, and does not come with a rising false-flag rate. Looking at these
metrics together prevents a misleading result, such as a ``perfect'' defense that
simply blocks everything.

\section{Expected Behaviour and Evaluation Plan}

\noindent\textit{All results in this section are \textbf{expected}. The system is
at the design stage and has not yet been built. No measured results are
reported here.}

\medskip
The system will be evaluated by running repeated attack--defend cycles on the
target model:

\begin{enumerate}
  \item \textbf{Round 1 (baseline):} send a set of attack prompts and measure how
  many get through. This gives the baseline \textbf{Attack Success Rate (ASR)}
  and the \textbf{Detection Rate}, or catch-rate, where
  $\mathrm{DR} = 1 - \mathrm{ASR}$.
  \item \textbf{Reflect and fix:} the reflection step finds why attacks were
  missed and proposes fixes. After human approval, the lessons are stored in
  memory and applied.
  \item \textbf{Round 2 and beyond:} attack again with the same prompts and with
  new ones, and measure the new catch-rate.
\end{enumerate}

\subsection{Expected Result}

The catch-rate is expected to rise across rounds as the system learns --- for
example, from a lower value in Round~1 to higher values in later rounds ---
without a human writing the fixes. Figure~\ref{fig:trend} shows this
\emph{expected and illustrative} trend. It is not a measured result.

\reportfigure[0.8\textwidth]{fig-6-1.png}{Expected (illustrative) improvement in defense across cycles. Illustrative only --- actual values will be measured during implementation}{fig:trend}

\subsection{Generalization Check}

A part of the attack prompts will be kept aside as a \emph{held-out set}. The
system will never use this set for building fixes. Testing on it later will show
whether the improvements are real, or whether the system has only memorized the
attacks it has already seen.

\subsection{Limitations}

\begin{itemize}
  \item The improvement curve will not be a perfect straight line. Real learning
  is bumpy, and some rounds may show little or no gain.
  \item Results will depend on the chosen model, the attack set, and the
  available compute.
  \item The actual numbers will be measured during implementation. This report
  presents the design and the evaluation plan only.
\end{itemize}
